%! Rational Functions and Vertical Asymptotes — Unit 3, Lesson 3.3 — Guided Notes (Answer Key)
\documentclass[10pt]{article}
\usepackage{precalculus-article}
\usepackage{precalculus-key}
\newcommand{\TallMath}[1]{$\displaystyle #1\rule[-1.4em]{0pt}{3.2em}$}

\begin{document}

\pageheader{Unit 3, Lesson 3.3}{Guided Notes --- Answer Key}
\namedateperiod

\vspace{0.1in}

\begin{objectivebox}
By the end of this lesson, I will be able to\ldots
\begin{enumerate}[itemsep=2pt]
  \item Identify the vertical asymptotes of a rational function by finding where the denominator
        equals zero and the numerator does \ans{not}. \hfill\textit{(LO 1.9.A)}
  \item Describe the behavior of a rational function near a vertical asymptote using
        \ans{one-sided} limit notation. \hfill\textit{(LO 1.9.A)}
\end{enumerate}
\end{objectivebox}

\vspace{0.1in}

\begin{vocabbox}
\termblanklong{Vertical asymptote}
\termblanklong{One-sided limit ($\lim_{x \to a^+}$)}
\termblanklong{One-sided limit ($\lim_{x \to a^-}$)}
\termblanklong{Multiplicity condition for a VA}
\end{vocabbox}

\vspace{0.1in}

\begin{hookbox}
Bacterial concentration near a drain: $C(x) = \dfrac{1}{x - 3}$, where $x$ is distance in cm.

\smallskip
\renewcommand{\arraystretch}{1.3}
\begin{tabular}{|c|c|c|c|c|c|c|}
\hline
$x$    & 2.9 & 2.99 & 2.999 & 3.001 & 3.01 & 3.1 \\
\hline
$C(x)$ & $-10$ & $-100$ & $-1{,}000$ & $1{,}000$ & $100$ & $10$ \\
\hline
\end{tabular}
\smallskip

As $x \to 3$, the output \ans{grows without bound (increases or decreases without limit)}.
We say $x = 3$ is a \ans{vertical asymptote}.
\end{hookbox}

\vspace{0.1in}

\begin{notesbox}{Vertical Asymptotes}
A rational function $r(x) = \dfrac{p(x)}{q(x)}$ has a \textbf{vertical asymptote at $x = a$} when:

\begin{itemize}
  \item $q(a) =$ \ans{$0$} \quad (the denominator is \ans{zero} at $a$), \textbf{AND}
  \item $p(a)$ \ans{$\ne$} $0$ \quad (the numerator is \textbf{not} zero at $a$).
\end{itemize}

\textit{Multiplicity rule (EK 1.9.A.1):} A VA also occurs at $x = a$ when the multiplicity of
$a$ in the \ans{denominator} is \emph{greater} than its multiplicity in the \ans{numerator}.

\medskip
\textbf{Example 1.} Find any vertical asymptotes of $r(x) = \dfrac{x + 2}{x - 3}$.

\begin{enumerate}[label=\textbf{Step \arabic*.}, itemsep=4pt]
  \item Set the denominator equal to zero and solve: \ans{$x - 3 = 0 \Rightarrow x = 3$}

  \item Check the numerator at each solution: at $x = 3$, numerator $=$ \ans{$5 \ne 0$}

  \item Conclusion: \ans{Vertical asymptote at $x = 3$.}
\end{enumerate}

\medskip
\textbf{Example 2.} Find any vertical asymptotes of $r(x) = \dfrac{x - 1}{(x-1)^2(x+2)}$.

\begin{enumerate}[label=\textbf{Step \arabic*.}, itemsep=4pt]
  \item Denominator zeros: \ans{$x = 1$ (mult.\ 2) and $x = -2$ (mult.\ 1)}

  \item At $x = 1$: numerator multiplicity \ans{1}, denominator multiplicity \ans{2}.
        Since \ans{2} $>$ \ans{1}: \ans{VA at $x = 1$}

  \item At $x = -2$: numerator value $=$ \ans{$-3 \ne 0$}. Since it is nonzero: \ans{VA at $x = -2$}
\end{enumerate}
\end{notesbox}

\vspace{0.1in}

\begin{notesbox}{One-Sided Limits Near a Vertical Asymptote}
For $r(x) = \dfrac{1}{x - 2}$, complete the table:

\smallskip
\renewcommand{\arraystretch}{1.3}
\begin{tabular}{|c|c|c|c|c|c|c|}
\hline
$x$    & 1.9 & 1.99 & 1.999 & 2.001 & 2.01 & 2.1 \\
\hline
$r(x)$ & \ans{$-10$} & \ans{$-100$} & \ans{$-1000$} & \ans{$1000$} & \ans{$100$} & \ans{$10$} \\
\hline
\end{tabular}
\smallskip

From the table, write the one-sided limits:

\[
  \lim_{x \to 2^-} r(x) = \ans{$-\infty$}
  \qquad\qquad
  \lim_{x \to 2^+} r(x) = \ans{$+\infty$}
\]

\textbf{Sign analysis.} When $x$ is slightly greater than 2, $(x - 2)$ is \ans{positive},
so $r(x) \to$ \ans{$+\infty$}.\\
When $x$ is slightly less than 2, $(x - 2)$ is \ans{negative},
so $r(x) \to$ \ans{$-\infty$}.
\end{notesbox}

\vspace{0.1in}

\begin{practicebox}
\begin{enumerate}[itemsep=10pt]
  \item For $r(x) = \dfrac{2x - 1}{x^2 - 5x + 6}$, find all vertical asymptotes.

        Factor the denominator: \ans{$(x-2)(x-3)$}

        Denominator zeros: \ans{$x = 2,\; x = 3$} \quad
        Check numerator at each: \ans{$2(2)-1=3\ne0$;\; $2(3)-1=5\ne0$}

        Vertical asymptote(s): \ans{$x = 2$ and $x = 3$}

  \item For $r(x) = \dfrac{1}{(x+3)^2}$, determine whether
        $\lim_{x\to -3^+} r(x) = +\infty$ or $-\infty$. Explain using sign analysis.

        \vspace{0.1in}
        \ans{$(x+3)^2 > 0$ for all $x \ne -3$, so $r(x) = 1/(x+3)^2 > 0$ near $x=-3$
        from both sides. Therefore $\lim_{x\to-3^+} r(x) = +\infty$.}
\end{enumerate}
\end{practicebox}

\end{document}
